\uuid{NHtJ}
\titre{Théorème de Rolle}
\auteur{bodin ; hanani ; mzari}
\organisation{exo7}
\licence{CC BY-NC-SA 4.0}
\datecreate{2020-01-20}
\niveau{L1}
\module{Analyse}
\chapitre{Dérivabilité des fonctions réelles}
\sousChapitre{Autre}
\difficulte{2}
\theme{Dérivées, théorème de Rolle}
\exo7qcmid{lille-1.derivees.019}

\contenu{
\begin{question}
\qclassification{124.00}
\qtags{motcle=théorème de Rolle}
On considère $\displaystyle f(x)=x^2-x$ sur l'intervalle $[0,1]$. Quelles sont les bonnes réponses ?
\begin{answers}
    \good{$f$ vérifie les hypothèses du théorème de Rolle et une valeur vérifiant la conclusion de ce théorème est $\displaystyle \frac{1}{2}$.}
    \bad{$f$ ne vérifie pas les hypothèses du théorème de Rolle.}
    \bad{$f$ ne vérifie pas les hypothèses du théorème des accroissements finis.}
    \good{$f$ vérifie les hypothèses du théorème des accroissements finis et une valeur vérifiant la conclusion de ce théorème est $\displaystyle \frac{1}{2}$.}
\end{answers}
\begin{explanations}
La fonction $f$ est continue sur $[0,1]$ et est dérivable sur $]0,1[$. Donc elle vérifie les hypothèses du théorème des accroissements finis, et, comme en plus $f(0)=0=f(1)$, elle vérifie aussi les hypothèses du théorème de Rolle. Les deux théorèmes impliquent l'existence de $c\in ]0,1[$ tel que $f'(c)=0$. Soit $\displaystyle c=\frac{1}{2}$.
\end{explanations}
\end{question}
}
